\originalsection{由偶极势恢复三个距离阶次}
将电偶极的势写为 $V=(4\pi\eps)^{-1}[a/r^2+b/(cr)]$, 其中 $a=\nn\cdot\pp(t_r)$、$b=\nn\cdot\dot\pp(t_r)$、$t_r=t-r/c$. 在源外 $r>0$, 有 $\nabla t_r=-\nn/c$, $\partial_jn_i=(\delta_{ij}-n_in_j)/r$. 对势中的两个项分别求导:
\begin{align*}
 \nabla(a/r^2)&=\frac{\pp-3\nn a}{r^3}-\frac{\nn b}{cr^2},\\
 \nabla[b/(cr)]&=\frac{\dot\pp-2\nn b}{cr^2}
 -\frac{\nn(\nn\cdot\ddot\pp)}{c^2r}.
\end{align*}
再减去 $\partial_t\vct A=\ddot\pp/(4\pi\eps c^2r)$, 第一行的 $r^{-3}$ 项、两行合并的 $r^{-2}$ 项、及最后的横向 $r^{-1}$ 项正好形成前面完整场式. 其中感应场也含三个法向投影系数, 因为梯度同时微分 $1/r$ 和单位方向; 仅微分时刻会漏掉它.

磁场由 $\nabla\times[\mu_0\dot\pp/(4\pi r)]$ 给出. 一项来自 $\nabla(1/r)=-\nn/r^2$, 另一项来自 $\nabla t_r$, 相加为
$\mu_0[\dot\pp\times\nn/r^2+\ddot\pp\times\nn/(cr)]/(4\pi)$.
这完成源外三个阶次的中文微分, 原点处仍须使用第4章分布而非普通函数公式. 变成磁偶极场时, 先限于源外无源真空再作对偶, 不能用源外变换猜接触项.
